% ============================================================
% Capítulo 03: Técnicas básicas de cuantificación de biomasa
% ============================================================

%\section{Técnicas de cuantificación de biomasa}
%cuenta en placa
%cuenta en cámara
%nefelometría
%gravimetría

% Consejo: Este capítulo cubre cómo medir la cantidad de microorganismos en un cultivo.

\section{Técnicas de cuantificación de biomasa}

La cuantificación de biomasa es esencial para monitorear el crecimiento microbiano y el progreso del bioproceso.

\subsection{Cuenta en placa}
Es una técnica de cuantificación de células viables (vivas).
\begin{enumerate}
    \item Se realiza una serie de diluciones del cultivo.
    \item Se siembra una alícuota de cada dilución en una placa con medio de cultivo sólido.
    \item Se incuba y se cuentan las colonias que crecen.
    \item El número de colonias se multiplica por el factor de dilución para obtener el número de unidades formadoras de colonias (UFC/mL).
\end{enumerate}
% Consejo: Es muy precisa para células viables, pero es lenta y solo cuenta células que pueden crecer en el medio.

\subsection{Cuenta en cámara}
Es una técnica de cuantificación de células totales (vivas y muertas).
\begin{enumerate}
    \item Se coloca una muestra del cultivo en una cámara de conteo (hemocitómetro o cámara de Petroff-Hausser) que tiene una cuadrícula de volumen conocido.
    \item Se observa al microscopio y se cuentan las células en varios cuadros.
    \item Se calcula el número de células/mL usando un factor de conversión.
\end{enumerate}
% Consejo: Es rápida y sencilla, pero no distingue entre células vivas y muertas y puede ser imprecisa.

\subsection{Nefelometría (Turbidimetría)}
Es una técnica de cuantificación de células totales basada en la dispersión de la luz.
\begin{enumerate}
    \item Se mide la turbidez o absorbancia del cultivo en un espectrofotómetro a una longitud de onda específica (generalmente 600 nm).
    \item La absorbancia es proporcional a la concentración de biomasa.
    \item Se requiere una curva de calibración (absorbancia vs. peso seco) para obtener valores de concentración (g/L).
\end{enumerate}
% Consejo: Es una técnica rápida, sencilla y no destructiva, pero la absorbancia puede verse afectada por otros componentes del medio.

\subsection{Gravimetría}
Es una técnica de cuantificación de células totales basada en el peso seco.
\begin{enumerate}
    \item Se toma un volumen conocido del cultivo y se centrifuga para obtener un pellet celular.
    \item Se lava el pellet para eliminar residuos del medio.
    \item Se seca el pellet en una estufa hasta peso constante.
    \item Se pesa el pellet y se calcula la concentración de biomasa (g/L de peso seco).
\end{enumerate}
% Consejo: Es directa y precisa, pero es destructiva, consume tiempo y no es adecuada para muestras con baja concentración de biomasa.